%% kapa: exam=1 round=33 year=2022 period=1 form=A subject=civil_law subject_ko=민법 seq=23 number=23
%% source: Q-Net 감정평가사 기출문제 (cst00309) artlSeq=5209997 file=raw/1008/5209997/2022년도 제33회 감정평가사 1차시험 1교시(민법,경제학,부동산학).pdf
%% conversion: kordoc 4.24.51 -> tools/gen1.py (auto); review=proofread-sonnet; flags=none
\begin{question}{23}
甲이 20년간 소유의 의사로 평온, 공연하게 乙소유의 X토지를 점유한 경우에 관한 설명으로 옳은 것을 모두 고른 것은? (다툼이 있으면 판례에 따름)
\begin{framedtext}
ㄱ. X토지가 미등기 상태라면 甲은 등기 없이도 X토지의 소유권을 취득한다. \par
ㄴ. 乙은 甲에 대하여 점유로 인한 부당이득반환청구를 할 수 있다. \par
ㄷ. 乙이 丙에게 X토지를 유효하게 명의신탁한 후 丙이 甲에 대해 소유자로서의 권리를 행사하는 경우, 특별한 사정이 없는 한 甲은 점유취득시효의 완성을 이유로 이를 저지할 수 있다.
\end{framedtext}
\begin{choices}
  \item ㄱ
  \item ㄷ
  \item ㄱ, ㄴ
  \item ㄴ, ㄷ
  \item ㄱ, ㄴ, ㄷ
\end{choices}
\end{question}
\answer{2}
